
En esta tarea se estudió un problema de selección de capítulos de anime bajo restricciones de minutos y energía, donde para cada anime solo es válido escoger un prefijo de capítulos. Se implementaron cuatro estrategias: fuerza bruta, programación dinámica y dos heurísticas greedy. El objetivo fue comparar correctitud, tiempo de ejecución, memoria utilizada y calidad de solución entre los distintos enfoques.

Para evaluar las implementaciones se generaron automáticamente casos de prueba pequeños y medianos, además de instancias controladas para aislar el efecto de la cantidad de animes, la cantidad total de capítulos, los minutos disponibles y la energía disponible. A partir de estas instancias se midieron tiempo y memoria, y se utilizó programación dinámica como referencia óptima para evaluar la calidad de los algoritmos greedy.

Los resultados muestran que la programación dinámica entrega la solución óptima, pero con mayor costo computacional; fuerza bruta solo es útil en casos pequeños; y las heurísticas greedy son mucho más rápidas y eficientes en memoria, aunque con una calidad inferior a la óptima. Entre ellas, la variante basada en energía tendió a acercarse más a la solución óptima que la basada en minutos.
